\documentclass[11pt,onecolumn]{IEEEtran}
\usepackage[a4paper]{geometry}
\usepackage{amsmath,amsfonts,amssymb}
\usepackage{color}
\usepackage{hyperref}

\topmargin      -12.0mm
\oddsidemargin  -7.0mm
\evensidemargin -7.0mm
\textheight     233.0mm
\textwidth      173.0mm
\columnsep        4.1mm
\parindent        1.5em
\headsep          6.3mm
\headheight        12pt
\lineskip           1pt
\normallineskip     1pt
\def\baselinestretch{1.1}

\begin{document}

\begin{center}
      {\bf\Large Response to the Reviewer Comments on Manuscript}\\[2mm]
      {\bf\Large ``Federated Wind Speed Forecasting in Finland: Geographic vs.\ Correlation Graphs with Synchronous FedGD''}\\[7mm]
      {\em\large corresponding author: Anonymous}\\[6.0mm]
      \today
\end{center}

\vspace{7mm}

We thank the reviewers and the teaching assistant for their constructive and detailed feedback.
We have addressed all raised issues in the revised manuscript.
\vspace{2mm}

Major modifications include:

\vspace{2mm}

\begin{itemize}
      \item The report title was modified to describe the technical approach of the project better. Previously it only mentioned the application domain, but I added the main idea (comparing graph structures) and the methods (FL with synchronous FedGD) to the title.
      \item All plot numbering from the figures removed and the latex auto-generated plot numbering is now used for referencing figures and plots.
      \item The feature count was corrected from 7 to 6, more mathematical notation was added where appropriate, and more justification for the use of the $\ell_2$ norm was added.
      \item The future work section was expanded from a short paragraph into a list of bullet points.
\end{itemize}

In the following, we respond to each comment point by point.
Section, figure, and equation numbers in the quoted comments refer to the original manuscript;
numbers in our responses refer to the revised manuscript.

\vspace{5mm}


%%%%%%%%%%%%%%%%%%%%%%%%%%%%%%%%%%%%%%%%%%%%
\begin{center}
      {\bf\large Comments of Reviewer \#1}
      \vspace{1mm}
\end{center}
%%%%%%%%%%%%%%%%%%%%%%%%%%%%%%%%%%%%%%%%%%%%

\begin{enumerate}

      \item[1.1] {\em The report should include the list of FMI station IDs or names used (e.g., as a table or
            supplementary reference), since the $\geq$80\% completeness filter applied to an undated API call
            cannot reliably reproduce the exact 50 stations.}

            \vspace*{2mm}

            \begin{itemize}
                  \item A complete station list was omitted given the 136 stations involved. Instead, the data query description in the numerical experiments section was strengthened with the exact API query, bounding box, variable names, and processing pipeline, so the filtering procedure can be reproduced exactly.
            \end{itemize}

            \vspace*{4mm}

      \item[1.2] {\em There is an inconsistency where figure numbers embedded inside plots (e.g., ``Figure 4'',
            ``Figure 7'', ``Figure 8'') do not match the caption labels below them (``Fig.\ 2'', ``Fig.\ 5'',
            ``Fig.\ 6''): these should be cleaned up before the final submission.}

            \vspace*{2mm}
            \begin{itemize}
                  \item All in-plot titles were removed. Now all figures only use the latex autogenerated figure numbering.
            \end{itemize}
            \vspace*{4mm}

      \item[1.3] {\em Some of the texts are overlapping and affect the readability of plots (Fig.\ 6 / Figure 8).
            Since there are already captions below each figure, the captions inside the plots seem unnecessary
            and reduce legibility.}

            \vspace*{2mm}
            \begin{itemize}
                  \item The plot text and titles were simplified so that the caption below each figure conveys the necessary information.
            \end{itemize}
            \vspace*{4mm}

\end{enumerate}


%%%%%%%%%%%%%%%%%%%%%%%%%%%%%%%%%%%%%%%%%%%%
\begin{center}
      {\bf\large Comments of Reviewer \#2}
      \vspace{1mm}
\end{center}
%%%%%%%%%%%%%%%%%%%%%%%%%%%%%%%%%%%%%%%%%%%%

\begin{enumerate}

      \item[2.1] {\em The report discusses the use of 7 features, but it only explains 6:
            3 lagged wind speeds (t-1, t-2, t-3), 1 lagged atmospheric pressure, 1 lagged temperature,
            1 lagged humidity.}

            \vspace*{2mm}
            \begin{itemize}
                  \item The feature count was corrected to 6 throughout the manuscript.
                        The weight vector $\mathbf{w}^{(i)} \in \mathbb{R}^7$ has 7 components because the
                        feature vector is augmented with a bias term, not because there are 7 features.
            \end{itemize}
            \vspace*{4mm}

      \item[2.2] {\em Feature normalization is repeated unnecessarily in the Problem Formulation section.}

            \vspace*{2mm}
            \begin{itemize}
                  \item The repeated normalization description was removed from the Problem Formulation section.
                        Feature and label normalization are now described only once, in Section 2.
            \end{itemize}
            \vspace*{4mm}

      \item[2.3] {\em GTVMin (and FedGD) is discussed in Section 2 before being explained in Section 3.
            It is also not cited. The reader should not be assumed to be previously familiar with GTVMin.}

            \vspace*{2mm}
            \begin{itemize}
                  \item A forward reference ``(defined in Section~3)'' and a citation
                        were added at the first mention of GTVMin in Section~2.
                        FedGD is also now cited when mentioned the first time.
            \end{itemize}
            \vspace*{4mm}

      \item[2.4] {\em For the Geographic k-NN system, the edge weight kernel explanation is a bit unclear:
            it uses a Gaussian kernel (not mentioned or cited); the report makes it sound like all of
            Equation~2 is the Haversine distance; the Haversine formula is not shown or cited.}

            \vspace*{2mm}
            \begin{itemize}
                  \item The edge weight formula (Eq.~2) is now described correctly as a Gaussian kernel.
                  \item The variable $d_{ij}$ in Eq.~2 is now clearly explained as the Haversine distance
                        between stations $i$ and $j$, so it is clear that the Haversine formula gives $d_{ij}$
                        and not the full equation.
                  \item A citation to the Haversine formula was added.
            \end{itemize}
            \vspace*{4mm}

      \item[2.5] {\em Some subsections would improve the report's readability.}

            \vspace*{2mm}
            \begin{itemize}
                  \item Subsections were not added. The spconf IEEE template is a well recognized
                        template and custom subsections might actually hurt readability if readers
                        are expecting a certain template structure.
            \end{itemize}
            \vspace*{4mm}

      \item[2.6] {\em The figures have inconsistent number labeling in their titles and in the report.}

            \vspace*{2mm}
            \begin{itemize}
                  \item All in-plot figure numbering removed.
            \end{itemize}
            \vspace*{4mm}

\end{enumerate}


%%%%%%%%%%%%%%%%%%%%%%%%%%%%%%%%%%%%%%%%%%%%
\begin{center}
      {\bf\large Comments of Teaching Assistant (Alexander Pavlyuk)}
      \vspace{1mm}
\end{center}
%%%%%%%%%%%%%%%%%%%%%%%%%%%%%%%%%%%%%%%%%%%%

\begin{enumerate}

      \item[TA.1] {\em Criterion 1: The title is not informative. Please indicate (i) the main idea (comparing
            graph structures), (ii) the application domain (wind prediction), and (iii) the methods
            (FL, Synchronous FedGD).}

            \vspace*{2mm}
            \begin{itemize}
                  \item The title was changed to ``Federated Wind Speed Forecasting in Finland:
                        Geographic vs.\ Correlation Graphs with Synchronous FedGD''.
                        It now conveys (i) the comparison of graph structures (Geographic vs.\ Correlation Graphs),
                        (ii) the application domain (wind speed forecasting in Finland), and
                        (iii) the method (Synchronous FedGD).
            \end{itemize}
            \vspace*{4mm}

      \item[TA.2] {\em Criterion 6: The label was not explicitly defined. Mathematical notation should be added
            to improve readability.}

            \vspace*{2mm}
            \begin{itemize}
                  \item The prediction label is now formally defined as $y^{(i)}_t = ws^{(i)}_{t+1}$
                        (next-hour wind speed at station $i$) at the start of Section~2.
                  \item The full feature vector is now written out as a math expression with all six
                        components named and their physical units stated.
            \end{itemize}
            \vspace*{4mm}

      \item[TA.3] {\em Criterion 9: More arguments in favor of the squared $\ell_2$ norm can be made.}

            \vspace*{2mm}
            \begin{itemize}
                  \item Two more arguments for the squared $\ell_2$ norm were added in Section 3:
                        \begin{enumerate}
                              \item it penalizes large parameter disagreements quadratically, encouraging
                                    more agreement between neighbors
                              \item adding a strictly convex penalty to the sum of convex local losses makes
                                    the full GTVMin objective strongly convex for $\alpha > 0$, guaranteeing a
                                    unique global minimizer
                        \end{enumerate}
            \end{itemize}
            \vspace*{4mm}

      \item[TA.4] {\em Criterion 14: Figure title/numbering issues; y-axis label out of bounds on Figure 3;
            in-plot title overlaps; station coordinates used instead of names.}

            \vspace*{2mm}
            \begin{itemize}
                  \item All in-plot figure numbers were removed and now the report only relies on latex automatic figure numbering
                  \item The y-axis label that was clipped out of bounds on Fig.~3 was moved to be in frame.
                  \item Station coordinates used to label the best and worst performing stations
                        were replaced with their FMI station names for readability.
            \end{itemize}
            \vspace*{4mm}

      \item[TA.5] {\em Criterion 16: I strongly suggest including more detailed potential improvements
            and presenting them in a list.}

            \vspace*{2mm}
            \begin{itemize}
                  \item The future directions were rewritten from a short paragraph into a bullet list.
                  \item The future directions contain reasons why the directions would strengthen the understanding of the application.
            \end{itemize}
            \vspace*{4mm}

\end{enumerate}

\vspace{5mm}

\end{document}
